\documentclass[10pt,a4paper]{amsart}
\usepackage{amsmath,amssymb,amsthm,graphicx}
% Retained minimal preamble for the Stacks Project chapter flat.tex (More on Flatness).
% Verbatim excerpt of the chapter preamble: Xy-pic setup (native preamble lines 14-18) and the native
% theorem declarations with the mathematical macros (native preamble lines 154-217). Copied mechanically;
% nothing added, removed or redefined. The wrapper class and the amsmath/amssymb/amsthm/graphicx base
% packages are supplied by the assembler, so the class/externaldocument part of the native preamble is omitted.
% For commutative diagrams we use Xy-pic
\usepackage[all]{xy}

% We use 2cell for 2-commutative diagrams.
\xyoption{2cell}
% Theorem environments.
%
\theoremstyle{plain}
\newtheorem{theorem}[subsection]{Theorem}
\newtheorem{proposition}[subsection]{Proposition}
\newtheorem{lemma}[subsection]{Lemma}

\theoremstyle{definition}
\newtheorem{definition}[subsection]{Definition}
\newtheorem{example}[subsection]{Example}
\newtheorem{exercise}[subsection]{Exercise}
\newtheorem{situation}[subsection]{Situation}

\theoremstyle{remark}
\newtheorem{remark}[subsection]{Remark}
\newtheorem{remarks}[subsection]{Remarks}

\numberwithin{equation}{subsection}

% Macros
%
\def\lim{\mathop{\mathrm{lim}}\nolimits}
\def\colim{\mathop{\mathrm{colim}}\nolimits}
\def\Spec{\mathop{\mathrm{Spec}}}
\def\Hom{\mathop{\mathrm{Hom}}\nolimits}
\def\Ext{\mathop{\mathrm{Ext}}\nolimits}
\def\SheafHom{\mathop{\mathcal{H}\!\mathit{om}}\nolimits}
\def\SheafExt{\mathop{\mathcal{E}\!\mathit{xt}}\nolimits}
\def\Sch{\mathit{Sch}}
\def\Mor{\mathop{\mathrm{Mor}}\nolimits}
\def\Ob{\mathop{\mathrm{Ob}}\nolimits}
\def\Sh{\mathop{\mathit{Sh}}\nolimits}
\def\NL{\mathop{N\!L}\nolimits}
\def\CH{\mathop{\mathrm{CH}}\nolimits}
\def\proetale{{pro\text{-}\acute{e}tale}}
\def\etale{{\acute{e}tale}}
\def\QCoh{\mathit{QCoh}}
\def\Ker{\mathop{\mathrm{Ker}}}
\def\Im{\mathop{\mathrm{Im}}}
\def\Coker{\mathop{\mathrm{Coker}}}
\def\Coim{\mathop{\mathrm{Coim}}}

% Boxtimes
%
\DeclareMathSymbol{\boxtimes}{\mathbin}{AMSa}{"02}

%
% Macros for moduli stacks/spaces
%
\def\QCohstack{\mathcal{QC}\!\mathit{oh}}
\def\Cohstack{\mathcal{C}\!\mathit{oh}}
\def\Spacesstack{\mathcal{S}\!\mathit{paces}}
\def\Quotfunctor{\mathrm{Quot}}
\def\Hilbfunctor{\mathrm{Hilb}}
\def\Curvesstack{\mathcal{C}\!\mathit{urves}}
\def\Polarizedstack{\mathcal{P}\!\mathit{olarized}}
\def\Complexesstack{\mathcal{C}\!\mathit{omplexes}}
% \Pic is the operator that assigns to X its picard group, usage \Pic(X)
% \Picardstack_{X/B} denotes the Picard stack of X over B
% \Picardfunctor_{X/B} denotes the Picard functor of X over B
\def\Pic{\mathop{\mathrm{Pic}}\nolimits}
\def\Picardstack{\mathcal{P}\!\mathit{ic}}
\def\Picardfunctor{\mathrm{Pic}}
\def\Deformationcategory{\mathcal{D}\!\mathit{ef}}

% --- appended 2026-10-12 (Chief): the project own definitions of the reference/slogan/history
% display environments, copied mechanically from the shared Stacks preamble lines 8-12
% (\usepackage{verbatim} plus the three \newenvironment lines, which the chapter preamble keeps
% verbatim from the project preamble). Nothing added, removed or redefined: these environments
% carry no mathematical content and the project itself typesets them as comments. Without these
% three lines a counted statement whose own text uses them fails with Environment slogan undefined.
\usepackage{verbatim}
\newenvironment{reference}{\comment}{\endcomment}
\newenvironment{slogan}{\comment}{\endcomment}
\newenvironment{history}{\comment}{\endcomment}
\usepackage{cleveref}
\crefname{originallemma}{original lemma}{original lemmas}
\Crefname{originallemma}{Original lemma}{Original lemmas}
\makeatletter
\newlabel{limits-lemma-approximate-given-relative-dimension}{{05M6}{source}{}{sourcepointer-1}{}}
\newlabel{limits-lemma-approximate-given-relative-dimension@cref}{{[originallemma][0][]05M6}{[1][0][]source}{}{}{}}
\makeatother
\crefname{originallemma}{original lemma}{original lemmas}
\Crefname{originallemma}{Original lemma}{Original lemmas}
\makeatletter
\newlabel{limits-lemma-descend-closed-immersion-finite-presentation}{{01ZP}{source}{}{sourcepointer-2}{}}
\newlabel{limits-lemma-descend-closed-immersion-finite-presentation@cref}{{[originallemma][0][]01ZP}{[1][0][]source}{}{}{}}
\makeatother
\crefname{originallemma}{original lemma}{original lemmas}
\Crefname{originallemma}{Original lemma}{Original lemmas}
\makeatletter
\newlabel{limits-lemma-descend-finite-presentation}{{01ZM}{source}{}{sourcepointer-3}{}}
\newlabel{limits-lemma-descend-finite-presentation@cref}{{[originallemma][0][]01ZM}{[1][0][]source}{}{}{}}
\makeatother
\crefname{originallemma}{original lemma}{original lemmas}
\Crefname{originallemma}{Original lemma}{Original lemmas}
\makeatletter
\newlabel{limits-lemma-descend-module-flat-finite-presentation}{{05LY}{source}{}{sourcepointer-4}{}}
\newlabel{limits-lemma-descend-module-flat-finite-presentation@cref}{{[originallemma][0][]05LY}{[1][0][]source}{}{}{}}
\makeatother
\crefname{originallemma}{original lemma}{original lemmas}
\Crefname{originallemma}{Original lemma}{Original lemmas}
\makeatletter
\newlabel{limits-lemma-limit-dimension}{{05M5}{source}{}{sourcepointer-5}{}}
\newlabel{limits-lemma-limit-dimension@cref}{{[originallemma][0][]05M5}{[1][0][]source}{}{}{}}
\makeatother
\crefname{originallemma}{original lemma}{original lemmas}
\Crefname{originallemma}{Original lemma}{Original lemmas}
\makeatletter
\newlabel{more-morphisms-lemma-flat-locus-base-change}{{047C}{source}{}{sourcepointer-6}{}}
\newlabel{more-morphisms-lemma-flat-locus-base-change@cref}{{[originallemma][0][]047C}{[1][0][]source}{}{}{}}
\makeatother
\crefname{originallemma}{original lemma}{original lemmas}
\Crefname{originallemma}{Original lemma}{Original lemmas}
\makeatletter
\newlabel{more-morphisms-theorem-openness-flatness}{{0399}{source}{}{sourcepointer-7}{}}
\newlabel{more-morphisms-theorem-openness-flatness@cref}{{[originallemma][0][]0399}{[1][0][]source}{}{}{}}
\makeatother
\crefname{originallemma}{original lemma}{original lemmas}
\Crefname{originallemma}{Original lemma}{Original lemmas}
\makeatletter
\newlabel{morphisms-lemma-dimension-fibre-after-base-change}{{02FY}{source}{}{sourcepointer-8}{}}
\newlabel{morphisms-lemma-dimension-fibre-after-base-change@cref}{{[originallemma][0][]02FY}{[1][0][]source}{}{}{}}
\makeatother
\crefname{originallemma}{original lemma}{original lemmas}
\Crefname{originallemma}{Original lemma}{Original lemmas}
\makeatletter
\newlabel{situation-flat-dimension-n}{{05MT}{source}{}{sourcepointer-9}{}}
\newlabel{situation-flat-dimension-n@cref}{{[originallemma][0][]05MT}{[1][0][]source}{}{}{}}
\makeatother
\title{The functor of points where a module is flat in dimensions at least n is an fpqc sheaf and limit preserving}
\author{The Stacks Project Authors}
\date{stacks-project:flat.tex:local-snapshot-20261011; GNU FDL 1.2 or later}
\begin{document}
\maketitle
\noindent\textit{Editorial note.} All delivered mathematical text is the author's own, copied from the pinned Stacks Project chapter file flat.tex; no statement or proof step was rewritten, reordered or removed.\par\medskip
\noindent\textit{Editorial note.} Preamble provenance: the chapter's own retained minimal preamble (sources/stacks-flat/retained-preamble.tex, lines 1-85) is a verbatim excerpt of the native Xy-pic setup and of the native theorem declarations and mathematical macros, followed by the project's own reference/slogan/history display-environment definitions copied mechanically from the shared preamble; the wrapper class and the amsmath/amssymb/amsthm/graphicx base packages are supplied by the assembler, and extra\_packages is left empty because everything needed is already in the retained preamble. No macro was invented or redefined.\par\medskip
\noindent\textit{Editorial note.} Same-chapter pointers are located at their \textbackslash{}label lines inside the pinned chapter file; cross-chapter pointers are located in the project's own tag list staged in the snapshot (sources/stacks-chow/tags.txt), at the verbatim lines carrying their canonical tags.\par\medskip
\noindent\textit{Allowed standard prerequisites (editorial).}\begin{itemize}
\item The situation defining the functor of points where the module is flat in dimensions at least n. Same chapter flat.tex, situation labelled situation-flat-dimension-n (statement, proof line 6338). Fixes what F\_n is; the proof has to show the sheaf property and the limit preserving property for that functor.
\item The flat locus and its behaviour under base change, and the dimension of fibres after base change. More on Morphisms, labels more-morphisms-lemma-flat-locus-base-change (proof line 6363) and more-morphisms-theorem-openness-flatness (proof lines 6399); Morphisms, label morphisms-lemma-dimension-fibre-after-base-change (proof line 6368). The first relates the non flat loci upstairs and downstairs so that flatness in dimensions at least n descends along a flat base change; openness of flatness identifies the open flat locus in the limit preserving part.
\item Approximation and descent lemmas from the theory of limits. Limits, labels limits-lemma-approximate-given-relative-dimension (proof line 6402), limits-lemma-descend-finite-presentation and limits-lemma-descend-closed-immersion-finite-presentation (proof lines 6406-6407), limits-lemma-limit-dimension (proof line 6410) and limits-lemma-descend-module-flat-finite-presentation (proof line 6412). Supply a finitely presented closed subscheme of dimension below n containing the non flat locus, descend it to a stage together with its dimension bound and with the flatness of the module off it, which yields flatness in dimensions at least n at that stage.
\end{itemize}\medskip
\ifdefined\hypertarget\hypertarget{sourcepointer-1}{}\fi
\noindent\textit{Source reference.} Original lemma 05M6: sources/stacks-chow/tags.txt, source lines 7260--7260. Navigation pointer for \textbackslash{}ref\{limits-lemma-approximate-given-relative-dimension\} used inside the counted span (line 6402); the argument that consumes it is retained in the span.\par\medskip
\ifdefined\hypertarget\hypertarget{sourcepointer-2}{}\fi
\noindent\textit{Source reference.} Original lemma 01ZP: sources/stacks-chow/tags.txt, source lines 2602--2602. Navigation pointer for \textbackslash{}ref\{limits-lemma-descend-closed-immersion-finite-presentation\} used inside the counted span (line 6407); the argument that consumes it is retained in the span.\par\medskip
\ifdefined\hypertarget\hypertarget{sourcepointer-3}{}\fi
\noindent\textit{Source reference.} Original lemma 01ZM: sources/stacks-chow/tags.txt, source lines 2599--2599. Navigation pointer for \textbackslash{}ref\{limits-lemma-descend-finite-presentation\} used inside the counted span (line 6406); the argument that consumes it is retained in the span.\par\medskip
\ifdefined\hypertarget\hypertarget{sourcepointer-4}{}\fi
\noindent\textit{Source reference.} Original lemma 05LY: sources/stacks-chow/tags.txt, source lines 7252--7252. Navigation pointer for \textbackslash{}ref\{limits-lemma-descend-module-flat-finite-presentation\} used inside the counted span (line 6412); the argument that consumes it is retained in the span.\par\medskip
\ifdefined\hypertarget\hypertarget{sourcepointer-5}{}\fi
\noindent\textit{Source reference.} Original lemma 05M5: sources/stacks-chow/tags.txt, source lines 7259--7259. Navigation pointer for \textbackslash{}ref\{limits-lemma-limit-dimension\} used inside the counted span (line 6410); the argument that consumes it is retained in the span.\par\medskip
\ifdefined\hypertarget\hypertarget{sourcepointer-6}{}\fi
\noindent\textit{Source reference.} Original lemma 047C: sources/stacks-chow/tags.txt, source lines 5503--5503. Navigation pointer for \textbackslash{}ref\{more-morphisms-lemma-flat-locus-base-change\} used inside the counted span (line 6363); the argument that consumes it is retained in the span.\par\medskip
\ifdefined\hypertarget\hypertarget{sourcepointer-7}{}\fi
\noindent\textit{Source reference.} Original lemma 0399: sources/stacks-chow/tags.txt, source lines 4264--4264. Navigation pointer for \textbackslash{}ref\{more-morphisms-theorem-openness-flatness\} used inside the counted span (line 6399); the argument that consumes it is retained in the span.\par\medskip
\ifdefined\hypertarget\hypertarget{sourcepointer-8}{}\fi
\noindent\textit{Source reference.} Original lemma 02FY: sources/stacks-chow/tags.txt, source lines 3209--3209. Navigation pointer for \textbackslash{}ref\{morphisms-lemma-dimension-fibre-after-base-change\} used inside the counted span (line 6368); the argument that consumes it is retained in the span.\par\medskip
\ifdefined\hypertarget\hypertarget{sourcepointer-9}{}\fi
\noindent\textit{Source reference.} Original lemma 05MT: sources/stacks-flat/pinned-original-source.tex, source lines 6311--6311. Navigation pointer for \textbackslash{}ref\{situation-flat-dimension-n\} used inside the counted span (line 6338); the argument that consumes it is retained in the span.\par\medskip
\begin{lemma}
\label{lemma-flat-dimension-n}
In Situation \ref{situation-flat-dimension-n}.
\begin{enumerate}
\item The functor $F_n$ satisfies the sheaf property for the fpqc topology.
\item If $f$ is quasi-compact and locally of finite presentation
and $\mathcal{F}$ is of finite presentation, then the functor $F_n$ is
limit preserving.
\end{enumerate}
\end{lemma}

\begin{proof}
Let $\{T_i \to T\}_{i \in I}$ be an fpqc covering of schemes over $S$.
Set $X_i = X_{T_i} = X \times_S T_i$ and denote $\mathcal{F}_i$ the
pullback of $\mathcal{F}$ to $X_i$. Assume that $\mathcal{F}_i$
is flat over $T_i$ in dimensions $\geq n$ for all $i$. Let $t \in T$.
Choose an index $i$ and a point $t_i \in T_i$ mapping to $t$.
Consider the cartesian diagram
$$
\xymatrix{
X_{\Spec(\mathcal{O}_{T, t})} \ar[d] &
X_{\Spec(\mathcal{O}_{T_i, t_i})} \ar[d] \ar[l] \\
\Spec(\mathcal{O}_{T, t}) &
\Spec(\mathcal{O}_{T_i, t_i}) \ar[l]
}
$$
As the lower horizontal morphism is flat we see from
More on Morphisms, Lemma \ref{more-morphisms-lemma-flat-locus-base-change}
that the set $Z_i \subset X_{t_i}$ where $\mathcal{F}_i$ is not flat
over $T_i$ and the set $Z \subset X_t$ where $\mathcal{F}_T$ is not flat
over $T$ are related by the rule $Z_i = Z_{\kappa(t_i)}$. Hence we see
that $\mathcal{F}_T$ is flat over $T$ in dimensions $\geq n$ by
Morphisms, Lemma \ref{morphisms-lemma-dimension-fibre-after-base-change}.

\medskip\noindent
Assume that $f$ is quasi-compact and locally of finite presentation and
that $\mathcal{F}$ is of finite presentation.
In this paragraph we first reduce the proof of (2) to the case where
$f$ is of finite presentation.
Let $T = \lim_{i \in I} T_i$ be a directed limit of
affine $S$-schemes and assume that $\mathcal{F}_T$ is flat in dimensions
$\geq n$. Set $X_i = X_{T_i} = X \times_S T_i$ and denote $\mathcal{F}_i$
the pullback of $\mathcal{F}$ to $X_i$. We have to show that
$\mathcal{F}_i$ is flat in dimensions $\geq n$ for some $i$.
Pick $i_0 \in I$ and replace $I$ by $\{i \mid i \geq i_0\}$.
Since $T_{i_0}$ is affine (hence quasi-compact) there exist finitely
many affine opens $W_j \subset S$, $j = 1, \ldots, m$ and an affine open
overing $T_{i_0} = \bigcup_{j = 1, \ldots, m} V_{j, i_0}$
such that $T_{i_0} \to S$ maps $V_{j, i_0}$ into $W_j$.
For $i \geq i_0$ denote $V_{j, i}$ the inverse image of $V_{j, i_0}$
in $T_i$. If we can show, for each $j$, that there exists an $i$ such that
$\mathcal{F}_{V_{j, i_0}}$ is flat in dimensions $\geq n$, then
we win. In this way we reduce to the case that $S$ is affine.
In this case $X$ is quasi-compact and we can choose a finite
affine open covering $X = W_1 \cup \ldots \cup W_m$. In this case
the result for $(X \to S, \mathcal{F})$ is equivalent to the result for
$(\coprod W_j, \coprod \mathcal{F}|_{W_j})$. Hence we may assume that
$f$ is of finite presentation.

\medskip\noindent
Assume $f$ is of finite presentation and $\mathcal{F}$ is of finite
presentation. Let $U \subset X_T$ denote the open subscheme of points where
$\mathcal{F}_T$ is flat over $T$, see
More on Morphisms, Theorem \ref{more-morphisms-theorem-openness-flatness}.
By assumption the dimension of every fibre of $Z = X_T \setminus U$ over
$T$ has dimension $< n$. By
Limits, Lemma \ref{limits-lemma-approximate-given-relative-dimension}
we can find a closed subscheme $Z \subset Z' \subset X_T$
such that $\dim(Z'_t) < n$ for all $t \in T$ and such that
$Z' \to X_T$ is of finite presentation. By
Limits, Lemmas \ref{limits-lemma-descend-finite-presentation} and
\ref{limits-lemma-descend-closed-immersion-finite-presentation}
there exists an $i \in I$ and a closed subscheme $Z'_i \subset X_i$
of finite presentation whose base change to $T$ is $Z'$. By
Limits, Lemma \ref{limits-lemma-limit-dimension}
we may assume all fibres of $Z'_i \to T_i$ have dimension $< n$. By
Limits, Lemma \ref{limits-lemma-descend-module-flat-finite-presentation}
we may assume that $\mathcal{F}_i|_{X_i \setminus T'_i}$
is flat over $T_i$. This implies that $\mathcal{F}_i$ is
flat in dimensions $\geq n$; here we use that $Z' \to X_T$ is of finite
presentation, and hence the complement $X_T \setminus Z'$ is quasi-compact!
Thus part (2) is proved and the proof of the lemma is complete.
\end{proof}

\end{document}
